\section{Uniform contact cancellation on an exact injection level}
\label{n:coherent-contact-section}

A tracked path that meets no other tracked path contributes nothing
to the alternating sum over retained coordinates.  We use this
cancellation to obtain a sparse-level contraction estimate, with
constants uniform in the number of tracked cards.  Combining it
with the lower-diagram bound of Proposition~\ref{n:strong-tail}
will give dimension-power decay when the tracked fraction is very
small.

The main issue is controlling the norm of the surviving kernel.
We estimate its square by a return experiment with common and
independent paths.  Conditional on the return randomness, a
weighted forest count bounds the probability that every tracked
path meets another.  This avoids an independence assumption on
the encounter edges.

Let $n=2^d$ with $d\ge1$, and let $T_n$ be the average operator of a
single directed sweep.  Write $X=[n]_{\ne}^k$ for the ordered injective
$k$-tuples and $Y=[n]^k$ for all ordered $k$-tuples.  Both spaces below
have counting measure.  Let
$J:\ell^2(X)\longrightarrow\ell^2(Y)$ be extension by zero; it is an
isometry.  A function $f:X\to\mathbb C$ is called harmonic if, for each
slot, its sum over that slot is zero when the other slots are fixed.
The sum uses the positions not occupied by the fixed slots.  Equivalently,
$Jf$ has zero sum over each unrestricted coordinate of $Y$, with the
other coordinates fixed.  Denote this subspace by $\mathcal H_k$.

\begin{proposition}[Uniform sparse contact estimate]
\label{n:coherent-contact}
There is an absolute constant $C_*>0$ such that for every dyadic
$n=2^d\ge2$, every $1\le k<n$, and every partition
$\lambda=(n-k,\beta)\vdash n$,
\begin{equation}\label{n:coherent-contact-sharp}
 \|T_n(\lambda)\|_{\op}^2
 \le \min\left\{1,\left(\frac{C_*dk}{n}\right)^{k/2}\right\}.
\end{equation}
For $k=1$ this operator is zero.  In particular, with an absolute
constant $C$ independent of $n,k,\lambda$,
\begin{equation}\label{rec:coherent-contact-bound}
 \|T_n(\lambda)\|_{\op}^2
 \le C^k(1+d)^{Ck}(k/n)^{k/2}.
\end{equation}
The same statements hold with the coordinate order reversed.
\end{proposition}

\begin{proof}
We first prove an operator bound on $\mathcal H_k$.  The representation
transfer at the end will explain why it applies to $V_\lambda$.

\medskip\noindent\textbf{Alternating coordinate kernels.}\enspace
For $A\subseteq[k]$, $x\in X$, and $y\in Y$, define
\[
 F_A(x,y)=n^{|A|}\Pr\{\pi x_i=y_i\text{ for every }i\in A\},
 \qquad P_A(x,y)=n^{-k}F_A(x,y),
\]
where $\pi$ is a sweep permutation.  The empty constraint has probability
one.  The matrix $P_A$ has rows indexed by $X$, columns indexed by $Y$,
and row sums one.  Its row law has the coordinates in $A$ follow one
common sweep, while the other coordinates are independent uniform
positions.  Let $K_X$ be the usual sweep transition kernel on $X$,
acting on functions by averaging their values at the output.  Then
$K_X=P_{[k]}J$.  If $A\ne[k]$ and $f\in\mathcal H_k$, summing over
any omitted output coordinate gives $P_AJf=0$.  Hence
\begin{equation}\label{n:contact-alternating-kernel}
 K_Xf=\sum_{A\subseteq[k]}(-1)^{k-|A|}P_AJf
 \qquad(f\in\mathcal H_k).
\end{equation}

A path through a directed sweep is determined by its two endpoints:
at each stage the coordinates already updated equal the corresponding
output coordinates, and those not yet updated equal the input
coordinates.  This defines a path even when the proposed paths cannot
be realized simultaneously.  For $x\in X,y\in Y$, make a graph on
$[k]$ by joining two labels when their prescribed paths enter the same
switch at some stage.  Switches are indexed by both stage and pair of
positions.  Let $E(x,y)$ be the event that this graph has no isolated
vertex.

Suppose label $i$ is isolated.  For $i\notin A$, the prescribed path
of $i$ uses $d$ switches disjoint from every other prescribed path.
Its probability is $2^{-d}=1/n$, independently of the prescriptions
for $A$.  If those prescriptions conflict, both probabilities are
zero.  Thus $F_{A\cup\{i\}}(x,y)=F_A(x,y)$ in every case, and the
alternating sum in~\eqref{n:contact-alternating-kernel} cancels pointwise.
Define the nonnegative rectangular matrices
\[
 M_A(x,y)=P_A(x,y)\mathbf1_{E(x,y)}.
\]
We have proved the exact identity
\[
 K_Xf=\sum_{A\subseteq[k]}(-1)^{k-|A|}M_AJf,
 \qquad
 \|K_X|_{\mathcal H_k}\|_{\op}
 \le\sum_{A\subseteq[k]}\|M_A\|_{\ell^2(Y)\to\ell^2(X)}.
\]
For $k=1$ the event $E$ is empty and the operator on $\mathcal H_1$
is zero.  Assume henceforth that $k\ge2$.

\medskip\noindent\textbf{The square of a retained kernel.}\enspace
Fix $A$ and a row $x\in X$.  After sampling $y$ according to $P_A(x,\cdot)$,
run a fresh reverse sweep for the labels in $A$, and let each other
label run its own independent reverse path.  Such a private path uses
independent fair stay-or-swap choices in the successive directions.
Its final position is uniform.  Because the forward-supported tuple
$y_A$ is injective, this defines a stochastic kernel on all return
tuples $z\in Y$:
\[
 Q_A(y,z)=n^{-|A^c|}
 \Pr\{\pi^{-1}y_i=z_i\text{ for every }i\in A\}.
\]
Here $\pi$ is a fresh forward sweep; its inverse has the reverse-sweep
law.  For $x'\in X$, the equality of the prescribed-image events gives
$Q_A(y,x')=P_A(x',y)$.  Thus the restriction of $Q_A$ to injective return
tuples is the adjoint kernel; the extension to all $z$ restores its
row sum to one.

Let $E_{\rm rev}(y,z)$ say that the prescribed reverse paths have no
isolated vertex.  Reversing paths preserves all switch encounters, so
$E(x',y)=E_{\rm rev}(y,x')$.  Dropping the first-leg encounter condition
and then the injectivity constraint on $x'$ gives
\begin{align}
 \sum_{x'\in X}(M_AM_A^*)(x,x')
 &\le \sum_{y\in Y}P_A(x,y)
       \sum_{z\in Y}Q_A(y,z)\mathbf1_{E_{\rm rev}(y,z)}.
 \label{n:contact-square-law}
\end{align}
Terms with $P_A(x,y)=0$ can be omitted.  The right side is the probability
that the mixed return paths have no isolated vertex.  We now bound it
uniformly in both $A$ and $x$.

\medskip\noindent\textbf{Product bounds for the midpoint occupation.}\enspace
Write $a=|A|$.  There is one common network carrying the $a$ labels in
$A$, and one private network for each of the $k-a$ other labels.
A vertex in the next construction records a network and an initial
position in that network.  Distinct networks give distinct vertices,
even when the physical position agrees.

In the common network the occupied midpoint set is the image of the
fixed $a$-set $\{x_i:i\in A\}$ under a sweep.  If its current occupancy
marginals are $p_v$, then for every position set $D$,
\begin{equation}\label{n:contact-occupation-products}
 \Pr\{D\text{ is occupied}\}\le\prod_{v\in D}p_v.
\end{equation}
This holds at the deterministic initial set.  A fair switch replaces
the two endpoint marginals by their mean.  For $D$ containing exactly
one endpoint, average the two old bounds; for $D$ containing both,
the probability is unchanged while the product of those two marginals
can only increase; the other sets are unaffected.  This proves
preservation under each switch.  At the end of a sweep each card is
uniform, so every common-network marginal is $a/n$.

Each private network has one uniform occupied vertex, independently
of the common network and of all other private networks.  Assign to
a typed vertex $v$ the weight
\[
 p_v=\begin{cases}
 a/n,&v\text{ belongs to the common network},\\
 1/n,&v\text{ belongs to a private network}.
 \end{cases}
\]
If $a=0$, omit the common network.  The random set $\mathcal O$ of
occupied typed vertices has exactly $k$ elements and, for every fixed
typed set $D$,
\begin{equation}\label{n:contact-typed-products}
 \Pr\{D\subseteq\mathcal O\}\le\prod_{v\in D}p_v,
 \qquad \sum_vp_v=k.
\end{equation}
A set requiring two different vertices of one private network has
probability zero, which is consistent with this bound.

\medskip\noindent\textbf{A graph determined by the return randomness.}\enspace
Condition only on the fresh randomness used for the return trip.
Sample the entire common switch wiring.  For each private label,
sample its $d$ independent stay-or-swap bits and use those same bits
for every possible starting position of that private network.  This
is a coupling of the possible starts: for each fixed start it has
exactly the private-path law already defined.  It changes no entry
of $Q_A$.  A fixed private bit at a stage applies either the identity
or the flip of that coordinate to all positions.  Consequently, in
every network, the map from a starting position to the position
entering any specified stage is a permutation.

Join two distinct typed vertices when their return paths enter the
same physical switch at some stage.  For a fixed vertex and a fixed
target network there are at most $2d$ neighbors: at a given stage
the switch has two positions, and each has exactly one preimage
under the target network's position map.  The weighted degree
therefore satisfies
\begin{equation}\label{n:contact-weighted-degree}
 \sum_{w:\,w\sim v}p_w\le\Delta,
 \qquad \Delta=2dk/n.
\end{equation}
The midpoint occupation is independent of the fresh return randomness,
so~\eqref{n:contact-typed-products} remains valid under this conditioning.
The encounter graph of the return paths is exactly the graph induced
by $\mathcal O$ in this typed graph.

\medskip\noindent\textbf{Counting forests with uniform constants.}\enspace
Every graph on $k$ vertices without isolated vertices contains a
spanning forest whose components all have at least two vertices.
Root each component and order the children at each vertex.  If the
forest has $j$ components, then $1\le j\le\lfloor k/2\rfloor$.
For fixed $k,j$, there are at most $4^k$ plane rooted forest shapes
with ordered components: adding one common root embeds this class
in the plane rooted trees with $k+1$ vertices, counted by the $k$th
Catalan number, which is at most $4^k$.

For a fixed forest shape $\mathcal F$, give an embedding $\phi$ into the
typed graph the weight $\prod_{u\in V(\mathcal F)}p_{\phi(u)}$.
Initially require injectivity and require that every forest edge maps
to a graph edge.  We may increase the sum by dropping injectivity,
retaining all factors with their multiplicities in this weight.
Sum the leaves first: for a fixed image of its parent, each leaf
contributes at most $\Delta$ by~\eqref{n:contact-weighted-degree}.
Delete that leaf and continue.  At the end, each of the $j$ roots
contributes $\sum_vp_v=k$.  The full sum is at most
$k^j\Delta^{k-j}$.

The component orderings can be divided out before dropping injectivity.
Indeed a fixed embedded forest with $j$ distinct components gives
$j!$ distinct shape--embedding pairs by ordering those components,
even if some component shapes coincide.  Together with
\eqref{n:contact-typed-products}, the union bound over witnessing
forests consequently gives
\begin{equation}\label{n:contact-weighted-forest-count}
 \Pr\{\mathcal O\text{ has no isolated vertex}\mid\text{return randomness}\}
 \le4^k\sum_{1\le j\le k/2}\frac{k^j}{j!}\Delta^{k-j}.
\end{equation}
This argument only uses a witnessing forest on the occupied set;
it does not assume that the encounter edges are independent.

If $\Delta\le1$, each $k-j\ge k/2$, and
\[
 4^k\sum_{1\le j\le k/2}\frac{k^j}{j!}\Delta^{k-j}
 \le4^k\Delta^{k/2}\sum_{j=0}^{\infty}\frac{k^j}{j!}
 =(4e)^k\Delta^{k/2}.
\]
This is uniform in the conditioned return graph and hence also holds
after averaging its randomness.  By~\eqref{n:contact-square-law}, every
row sum of the symmetric nonnegative matrix $M_AM_A^*$ has that upper
bound.  Schur's test, or its maximum-row-sum bound, gives
\[
 \|M_A\|_{\op}^2\le(4e)^k\Delta^{k/2}.
\]
There are $2^k$ choices of $A$, so the triangle inequality in
\eqref{n:contact-alternating-kernel} yields
\begin{equation}\label{n:contact-harmonic-bound}
 \|K_X|_{\mathcal H_k}\|_{\op}^2
 \le(16e)^k(2dk/n)^{k/2}.
\end{equation}
For $\Delta>1$, the contraction bound $\|K_X\|_{\op}\le1$ gives the
same conclusion after enlarging the absolute constant.  For example,
$C_*=512e^2$ suffices in~\eqref{n:coherent-contact-sharp}.

\medskip\noindent\textbf{Passing to the exact representation level.}\enspace
The permutation representation on $X$ is induced from the trivial
representation of the stabilizer $S_{n-k}$.  By Young branching and
Frobenius reciprocity, $V_\lambda$, with $\lambda_1=n-k$, occurs in
this representation: remove the $k$ boxes below its first row in a
valid corner-removal order.  It does not occur on $(k-1)$-tuples,
because occurrence there would require a first row of length at
least $n-k+1$.

For any specified slot, the functions lifted from the other $k-1$
slots form an invariant subspace containing no copy of $V_\lambda$.
Its orthogonal complement consists exactly of functions whose sum
in that slot is zero.  Thus every $V_\lambda$-isotypic component on
$X$ lies in $\mathcal H_k$.  The averaging operator $K_X$ on functions
has the Fourier action of $T_n^*$ under the convention that a
permutation acts by pullback through its inverse.  Its norm on a
copy of $V_\lambda$ is therefore $\|T_n(\lambda)\|_{\op}$.
Applying~\eqref{n:contact-harmonic-bound} proves the proposition.
The proof uses only that each coordinate is updated once and works
unchanged in the reversed coordinate order.
\end{proof}

\begin{remark}[Counting roots among physical positions]
For comparison with counting roots among physical positions, the
right side of~\eqref{n:contact-weighted-forest-count} equals
\[
 4^k(k/n)^k
 \sum_{1\le j\le k/2}\frac{n^j}{j!}(2d)^{k-j}.
\]
For $j\le n/2$, the bound
\[
 \log\frac{n^j}{(n)_j}
 =\sum_{r=0}^{j-1}-\log(1-r/n)
 \le\sum_{r=0}^{j-1}\frac{r}{n-r}\le j
\]
shows that $n^j/j!\le e^j\binom nj$.  Thus we also obtain the
root-count bound
\begin{equation}\label{rec:coherent-forest-count}
 \Pr\{\mathcal O\text{ has no isolated vertex}\mid\text{return randomness}\}
 \le(k/n)^k\sum_{1\le j\le k/2}(4e)^k\binom nj(2d)^{k-j}.
\end{equation}
The weighted version identifies all types and occupation factors
explicitly; both forms retain the same number $k-j$ of contact costs.

\end{remark}

\begin{corollary}[Exponential sparse-level saving]
\label{n:coherent-contact-sparse-saving}
There are absolute constants $C,d_0>0$ such that for every dyadic
$n=2^d$ and $\lambda=(n-k,\beta)\vdash n$ satisfying
$1\le k<n$ and $\log(n/k)\ge d^{3/4}$,
\[
 \|T_n(\lambda)\|_{\op}^2
 \le e^{Ck}(k/n)^{k/4}.
\]
For $d\ge d_0$, the factor $e^{Ck}$ can be omitted.
\end{corollary}
\begin{proof}
Put $s=\log(n/k)$.  For a nonzero operator,
\eqref{n:coherent-contact-sharp} bounds the logarithm of its squared norm by
$-ks/2+(k/2)\log(C_*d)$.  Since $\log(C_*d)=o(d^{3/4})$, it is at
most $-ks/4$ for all sufficiently large $d$ in the stated range.
The remaining finitely many $d$ are covered by increasing $C$, using
$\|T_n\|_{\op}\le1$ and $s\le d\log2$.  A zero operator is covered
directly.
\end{proof}

\begin{corollary}\label{n:coherent-sparse-consequence}
There are absolute $c>0$ and $d_0$ such that, for every
$N=2^d$ with $d\ge d_0$ and every partition
$\lambda=(N-k,\beta)\vdash N$ satisfying
$1\le k<N$ and $\log(N/k)\ge d^{3/4}$,
\[
 \|T_N(\lambda)\|_{\mathrm{op}}^2\le D_\lambda^{-c}.
\]
\end{corollary}
\begin{proof}
Put $s=\log(N/k)$ and $L=\log D_\beta$.  For sufficiently large
$d$, \eqref{rec:coherent-contact-bound} has logarithm at most
$-ks/4$, since its positive error is $O(k\log(1+d))$.
Proposition~\ref{n:strong-tail} gives the second logarithmic bound
$C_1k-aL$.  If $L\le(2C_1/a)k$, then
\[
 \log D_\lambda\le k(1+s)+L\le C_2ks,
\]
so the first bound suffices.  If $L>(2C_1/a)k$, the second bound is
at most $-aL/2$.  Taking the better of the two bounds then gives
\[
 \log\|T_N(\lambda)\|_{\mathrm{op}}^2
 \le-\max\{ks/4,aL/2\}
 \le-c_1(ks+L).
\]
The tableau inequality
$D_\lambda\le\binom NkD_\beta\le(eN/k)^kD_\beta$
again proves the result.  Zero operator norm is covered directly.
\end{proof}
